\section*{الفصل \ref{ch:b1:intermolecular-forces-solvents} --- القوى بين الجزيئات والمذيبات}

\begin{solution}{exo:b1:intermolecular-forces-solvents:1}
الأرغون: لندن وحده. كلوريد الهيدروجين: كيسوم وديباي ولندن، والغلبة
للندن (\cref{exo:b1:intermolecular-forces-solvents:8})؛ والكلور
أكبر من أن يكوّن \omterm{def:b1:intermolecular-forces-solvents:hbond}{روابط هيدروجينية} قوية. الماء: الغلبة \omterm{def:b1:intermolecular-forces-solvents:hbond}{للروابط الهيدروجينية}، مع
حدود فان دير فالس الثلاثة. الميثان: لندن وحده. ثنائي اليود: لندن
وحده، وهو قوي لأن الجزيء كبير وشديد القابلية للاستقطاب.
\end{solution}

\begin{solution}{exo:b1:intermolecular-forces-solvents:2}
بين جزيئاتها نفسها: \ce{CH3OH} (\ce{O-H})، و\ce{CH3NH2}
(\ce{N-H}) و\ce{HF}. أما \ce{CH3OCH3} و\ce{CH3F} فلهما \omterm{def:b1:lewis-resonance-vsepr:lewis}{أزواج حرة} على O
أو F لكن لا هيدروجين عليهما: فلا يستطيعان إلا أن يستقبلا \omterm{def:b1:intermolecular-forces-solvents:hbond}{روابط هيدروجينية}، من
الماء مثلًا. و\ce{CH4} لا يمنح ولا يستقبل.
\end{solution}

\begin{solution}{exo:b1:intermolecular-forces-solvents:3}
الماء والإيثانول: قطبيان، بروتونيان. البروبانون والأسيتونتريل: قطبيان،
لابروتونيان. ثنائي كلورو الميثان: ضعيف القطبية، لابروتوني. الهكسان: غير قطبي، لابروتوني.
\end{solution}

\begin{solution}{exo:b1:intermolecular-forces-solvents:4}
ثنائي اليود غير قطبي: ففي الهكسان يستبدل تآثرات لندن بتآثرات
لندن، بينما في الماء كان سيكسر \omterm{def:b1:intermolecular-forces-solvents:hbond}{روابط هيدروجينية} دون أن
يعوضها. وكلوريد الصوديوم أيوني: فلا يستطيع فصل أيوناته إلا مذيب عالي
السماحية، ولا يستطيع ذوبنتها إلا \omterm{def:b1:intermolecular-forces-solvents:solvent-classes}{مذيب قطبي}؛ والهكسان
($\varepsilon_r = 1.89$) لا يفعل هذا ولا ذاك.
\end{solution}

\begin{solution}{exo:b1:intermolecular-forces-solvents:5}
ليس للميثوكسي ميثان رابطة \ce{O-H}: فلا يستطيع أن يمنح \omterm{def:b1:intermolecular-forces-solvents:hbond}{روابط هيدروجينية} ويغلي
عند أدنى درجة. والميثانول والإيثانول كلاهما مترابط هيدروجينيًا؛ وللإيثانول، بمجموعة
\ce{CH2} زائدة، تآثرات لندن أقوى. ويكوّن حمض الإيثانويك
مثنويات مترابطة هيدروجينيًا وهو الأثقل: فيغلي عند أعلى درجة.
\end{solution}

\begin{solution}{exo:b1:intermolecular-forces-solvents:6}
في الفراغ $F = e^2/(4\pi\varepsilon_0 r^2) = (1.602 \times
10^{-19})^2/(4\pi \times 8.854 \times 10^{-12} \times (5.00 \times
10^{-10})^2) = \qty{9.23e-10}{N}$. وفي الماء $F/78.4 = \qty{1.18e-11}{N}$؛
وفي الهكسان $F/1.89 = \qty{4.88e-10}{N}$. فالماء يُضعف التجاذب بمقدار 41
مرة أكثر مما يضعفه الهكسان: فهو يستطيع إبقاء الأيونات متباعدة، والهكسان لا يستطيع.
\end{solution}

\begin{solution}{exo:b1:intermolecular-forces-solvents:7}
يكسر خلط البروبانون بالماء بعض \omterm{def:b1:intermolecular-forces-solvents:hbond}{الروابط الهيدروجينية} ماء--ماء لكنه
يكوّن روابط جديدة بين الماء وأكسجين \ce{C=O}، إضافة إلى
تآثرات ثنائي قطب--ثنائي قطب: فالحصيلة مواتية. أما الهكسان فلا يقدم
للماء إلا تآثرات لندن ضعيفة مقابل \omterm{def:b1:intermolecular-forces-solvents:hbond}{الروابط الهيدروجينية} التي
سيكسرها: فيبقى السائلان منفصلين.
\end{solution}

\begin{solution}{exo:b1:intermolecular-forces-solvents:8}
لجزيء \ce{HBr} إلكترونات أكثر وسحابة أكبر وأكثر قابلية للاستقطاب: \omterm{def:b1:intermolecular-forces-solvents:vdw}{فتآثر
لندن} فيه يفوق نظيره في \ce{HCl} بأكثر مما ينقص \omterm{def:b1:intermolecular-forces-solvents:vdw}{تآثر كيسوم}
فيه عنه. فالغلبة للندن.
\end{solution}

\begin{solution}{exo:b1:intermolecular-forces-solvents:9}
\babelsublr{\chemfig{CH_3-C(=[:60]O)-[:-60]O-[:0]H}} مقابل شريكه، ورابطتان
$\ce{O-H}\cdots\ce{O=C}$ تُغلقان حلقة ثمانية الذرات. في البنزين،
غير القطبي واللابروتوني، لا شيء ينافس هاتين الرابطتين: فيكون الحمض
مثنويًا ويسلك سلوك جسيم كتلته نحو \qty{120}{g/mol}. أما في الماء،
فكل جزيء من الحمض يترابط هيدروجينيًا مع الماء بدلًا من ذلك: فتتفكك المثنويات.
\end{solution}

\begin{solution}{exo:b1:intermolecular-forces-solvents:10}
كلوريد الصوديوم مكوّن من أيونات: والماء، العالي السماحية، يفصلها
(التفكك) ويُميهها (الأكسجين نحو \ce{Na+}، والهيدروجين
نحو \ce{Cl-}). وكلوريد الهيدروجين \omterm{def:b1:lewis-resonance-vsepr:dipole}{جزيء قطبي}: والماء، القطبي،
يحوّل الرابطة \ce{H-Cl} إلى زوج أيونات \ce{H3O+ Cl-} (التأين)،
ويفصل الزوج (التفكك) ويُميه الأيونات. أما الهكسان فغير قطبي
ومنخفض السماحية: فلا يؤيّن ولا يفكك، فلا أيونات
تحمل التيار.
\end{solution}

\begin{solution}{exo:b1:intermolecular-forces-solvents:11}
$(174.1 - 36.1)/5 = \qty{27.6}{\celsius}$ لكل \ce{CH2}. والزيادة
الأخيرة (من النونان إلى الديكان) \qty{23.4}{\celsius}، فينبغي أن يغلي الأونديكان
قرب $174 + 22 \approx \qty{196}{\celsius}$. فكل \ce{CH2} تضيف
الإسهام نفسه من \omterm{def:b1:intermolecular-forces-solvents:vdw}{تآثر لندن}، لكنه جزء يصغر شيئًا فشيئًا من مجموع
تآثرات الجزيء، في حين أن درجة الحرارة اللازمة للتغلب على
التآثرات صارت مرتفعة أصلًا.
% ledger: bp:C9H20
\end{solution}

\begin{solution}{exo:b1:intermolecular-forces-solvents:12}
الجزء \omterm{def:b1:intermolecular-forces-solvents:amphiphile}{المحب للماء}: رأس الكبريتات \ce{-OSO3^-} مع أيونه المرافق \ce{Na+}؛
والجزء \omterm{def:b1:intermolecular-forces-solvents:amphiphile}{الكاره للماء}: السلسلة ذات الاثنتي عشرة ذرة كربون. عند السطح، الرؤوس في
الماء والذيول في الهواء؛ وفي المذيلة، الذيول في الداخل والرؤوس في الخارج.
تذوب الذيول في دهن البقعة، وتبقيه الرؤوس على
تماس مع الماء، فيُرفع الدهن في مذيلات.
\end{solution}

\begin{solution}{pb:b1:intermolecular-forces-solvents:1}
\textbf{1.} \omterm{def:b1:intermolecular-forces-solvents:vdw}{تآثر لندن}، وهو الوحيد بين ذرات غير قطبية.
\textbf{2.} من الهيليوم إلى الزينون يكبر عدد الإلكترونات والحجم،
فتكبر \omterm{def:b1:periodicity:polarisability}{قابلية الاستقطاب} وتجاذب لندن.
\textbf{3.} للهيليوم إلكترونان تمسكهما نواته بقوة: فهو
أقل الذرات قابلية للاستقطاب، وأضعفها تجاذب لندن.
\textbf{4.} $68.8 - 36.1 = \qty{32.7}{\celsius}$: تضيف مجموعة \ce{CH2} واحدة
إلكترونات وسطح تماس، ومن ثم تجاذب لندن.
\textbf{5.} لندن، الذي يكبر مع \omterm{def:b1:periodicity:polarisability}{قابلية الاستقطاب}: $\ce{CH4} < \ce{SiH4} <
\ce{GeH4}$.
\textbf{6.} ليس الكربون كهرسلبيًا بما يكفي لاستقطاب \ce{C-H}
بقوة، وليس للميثان \omterm{def:b1:lewis-resonance-vsepr:lewis}{زوج حر} يستقبل به \omterm{def:b1:intermolecular-forces-solvents:hbond}{رابطة هيدروجينية}.
\textbf{7.} الماء: قطبي، بروتوني، وأعلاها سماحية؛ فهو
يفصل الأيونين كليهما ويذوبنهما.
\textbf{8.} $78.4/1.89 = 41$: التجاذب أضعف بمقدار 41 مرة في الماء.
\textbf{9.} الأسيتونتريل (أو البروبانون): قطبي بما يكفي لإذابة الملح،
ولابروتوني فيبقى الأنيون حرًا.
\textbf{10.} الهكسان: غير \omterm{def:b1:intermolecular-forces-solvents:miscibility}{قابل للامتزاج} بالماء، وغير قطبي مثل ثنائي اليود. فينتقل اليود
إلى الهكسان، الطبقة العليا (الهكسان أقل كثافة من الماء).
\textbf{11.} الإيثانول \omterm{def:b1:intermolecular-forces-solvents:miscibility}{قابل للامتزاج} بالماء: فلا تتكوّن طبقة ثانية.
\textbf{12.} يستقبل أكسجينه \omterm{def:b1:intermolecular-forces-solvents:hbond}{روابط هيدروجينية} من الماء، لكنه لا يستطيع
أن يمنح أيًّا منها، ومجموعتا الإيثيل فيه كارهتان للماء: فلا يذوب إلا
قليلًا.
\textbf{13.} $-66.4 - (-85.0) = \qty{18.6}{\celsius}$ لكل \omterm{def:b1:periodicity:table}{دورة}؛
و\ce{HF} المستكمَل: $-85.0 - 18.6 = \qty{-103.6}{\celsius}$.
\textbf{14.} يغلي \ce{HF} عند \qty{19.6}{\celsius}، أي \qty{123}{\celsius}
فوق الاستكمال.
\textbf{15.} $-62.5 - (-87.8) = \qty{25.3}{\celsius}$؛ و\ce{NH3}
المستكمَل: \qty{-113.1}{\celsius}؛ والفعلي $-33.4$، أي أعلى بمقدار \qty{80}{\celsius}.
\textbf{16.} $-88.1 - (-112.0) = \qty{23.9}{\celsius}$؛ و\ce{CH4}
المستكمَل: \qty{-135.9}{\celsius}؛ والفعلي $-161.5$، أي تحته. فالميثان لا يكوّن
\omterm{def:b1:intermolecular-forces-solvents:hbond}{روابط هيدروجينية}؛ إنه ببساطة صغير وضعيف القابلية للاستقطاب.
\textbf{17.} \babelsublr{\foreignlanguage{arabic}{الماء (نحو \qty{179}{\celsius})} $>$ \ce{HF} (123) $>$
\ce{NH3} (80)}.
\textbf{18.} يستطيع \ce{NH3} أن يمنح ثلاثًا لكنه لا يستقبل إلا واحدة (\omterm{def:b1:lewis-resonance-vsepr:lewis}{زوج حر} واحد)؛
ويستطيع \ce{HF} أن يستقبل ثلاثًا لكنه لا يمنح إلا واحدة. أما الماء فيمنح اثنتين ويستقبل
اثنتين: فكل جزيء يستطيع المشاركة في أربع \omterm{def:b1:intermolecular-forces-solvents:hbond}{روابط هيدروجينية}، فتتكوّن شبكة
ثلاثية الأبعاد كاملة.
\textbf{19.} لا يكوّن \ce{H2S} و\ce{H2Se} \omterm{def:b1:intermolecular-forces-solvents:hbond}{روابط هيدروجينية} ذات شأن:
فدرجتا غليانهما تتبعان منحى لندن الذي كان الماء سيتبعه
لولاها. أما الماء نفسه فهو الشذوذ المراد قياسه.
\textbf{20.} $-41.3 - (-60.3) = \qty{19.0}{\celsius}$ لكل \omterm{def:b1:periodicity:table}{دورة}.
\textbf{21.} $T(p) = -60.3 + 19.0\,(p - 3)$ بالوحدة \unit{\celsius}.
\textbf{22.} $T(2) = -60.3 - 19.0 = \qty{-79.3}{\celsius}$.
\textbf{23.} $100.0 - (-79.3) = \qty{179}{\celsius}$.
\textbf{24.} لولا روابطه الهيدروجينية لغلى الماء قرب
\textbf{\qty{-79}{\celsius}}، ولكان غازًا على الأرض.
\end{solution}
